\chapter{Lernziele Kompression}
\begin{todolist}
\item Ich kann Beispiele aus der Geschichte angeben, die aufzeigen, dass Kompression schon früh ein Bedürfnis war.
\item Ich kann den Unterschied zwischen verlustbehafteten und verlustfreien Komprimierungen erklären und Beispiele dafür angeben.
\item Ich weiss, was eine präfixfreie Kodierung ist, und wieso diese Eigenschaft wichtig ist.
\item Ich kann mit einem Baumdiagramm eine präfixfreie Kodierung finden.
\item Ich kann die absoluten und relativen Häufigkeiten von Zeichen in einem (kurzen) Text bestimmen.
\item Ich kann eine präfixfreie Kodierung mit der Methode der maximalen Balance für einen Text finden, von dem die relativen Häufigkeiten der Zeichen gegeben ist.
\item Ich kann ein Beispiel für eine relative Häufigkeit von Buchstaben angeben, bei welchem die Kodierung der maximalen Balance einen selteneren Buchstaben kürzer kodiert als einen häufigeren.
\item Ich kann eine Huffman-Kodierung für einen Text finden, von dem die relativen Häufigkeiten der Zeichen gegeben ist.
\item Ich kann ausgehend von der Kodierung und der absoluten Häufigkeiten der vorkommenden Zeichen berechnen, wie viele Bits für die Speicherung eines Textes nötig sind.
\item Ich kann bei gegebener relativer Häufigkeit der Zeichen einen (kurzen) Text arithmetisch Kodieren.
\item Ich kann bei gegebener relativer Häufigkeit der Zeichen eine arithmetische Kodierung dekomprimieren.
\end{todolist}